\section*{Chapter \ref{ch:rc:sabr-in-rates-and-the-volatility-cube} --- SABR in Rates and the Volatility Cube}

\begin{solution}{exo:rc:sabr-in-rates-and-the-volatility-cube:1}
No: its formula needs $F>0$ and $K>0$ (it takes $\ln(F/K)$ and powers of $FK$).
Any shift above 0.50\% makes $K+\zeta>0$ and $F+\zeta>0$; in practice one leaves a
margin, since the density near the floor misbehaves.
\end{solution}

\begin{solution}{exo:rc:sabr-in-rates-and-the-volatility-cube:2}
The at-the-money normal volatility is unchanged ($\alpha$ times the same factor):
the backbone is flat. Yesterday's at-the-money strike is now 20 basis points
below the forward, so its volatility is the smile's value at $F-K=+20$ basis
points: the whole smile has moved with the forward.
\end{solution}

\begin{solution}{exo:rc:sabr-in-rates-and-the-volatility-cube:3}
Within one section $\beta$ and $\rho$ both tilt the smile and are nearly
interchangeable, so a fit of both is unstable from day to day; $\beta$ is chosen
for the backbone the desk believes (how the at-the-money volatility moves with
the forward) and held fixed for the currency.
\end{solution}

\begin{solution}{exo:rc:sabr-in-rates-and-the-volatility-cube:4}
$25\times\bigl(1+\tfrac{2-3\times0.01}{24}\times0.36\times1\bigr) = 25\times1.02955 =
25.74$ basis points, the quoted at-the-money volatility.
\end{solution}

\begin{solution}{exo:rc:sabr-in-rates-and-the-volatility-cube:5}
The quotes were generated by a normal model; a \omterm{def:rc:sabr-in-rates-and-the-volatility-cube:shifted}{shifted lognormal model} with a
small shift has a strongly level-dependent volatility, since the shifted
forward $0.70\%$ is small, and its smile shape (curvature against skew) differs
from the normal one, so three parameters cannot match six quotes exactly. A
larger shift makes \omterm{def:rc:sabr-in-rates-and-the-volatility-cube:shifted}{shifted SABR} behave more like a normal model, and both the
fit and the wing agree better.
\end{solution}

\begin{solution}{exo:rc:sabr-in-rates-and-the-volatility-cube:6}
Buy the model's butterfly where the density is negative (long two receivers at
the middle strike, short one each at strikes a little above and below, around
0\%): the model prices it negative, so a desk marking with the 1\% fit would
pay to sell it. In practice those strikes are not quoted, bid--offer is wide,
and the market's own prices of such butterflies are positive: the arbitrage is
in the model, and a desk using it would be picked off.
\end{solution}

\begin{solution}{exo:rc:sabr-in-rates-and-the-volatility-cube:7}
0.90\%: at 0.85\% the calibrated smile's density is negative just above the
floor ($-0.80\%$ to $-0.79\%$), at 0.90\% it is positive from $-0.80\%$ to $+1\%$.
(Any shift must exceed 0.80\% for the quoted strikes to exist at all.)
\end{solution}

\begin{solution}{exo:rc:sabr-in-rates-and-the-volatility-cube:8}
The sections have different forwards, so the same absolute strike sits at
different moneyness in each; interpolating volatilities strike by strike mixes
at-the-money and wing points, and the result need not be a SABR smile or free
of arbitrage. Interpolate parameters (or moneyness-indexed volatilities), take
the forward from the curve, and check the density.
\end{solution}

\begin{solution}{pb:rc:sabr-in-rates-and-the-volatility-cube:1}
\textbf{1.} 60 basis points below the forward, 10 below the lowest quote
($-0.80\%$).
\textbf{2.} Its formula needs positive forwards and strikes.
\textbf{3.} \omterm{def:rc:sabr-in-rates-and-the-volatility-cube:shifted}{Shifted SABR}, which assumes the rate stays above $-\zeta$; \omterm{def:rc:sabr-in-rates-and-the-volatility-cube:normal}{normal SABR},
which lets the rate take any value, with normal tails.
\textbf{4.} 0.77 basis points (1\% shift), 0.11 (3\%), zero for \omterm{def:rc:sabr-in-rates-and-the-volatility-cube:normal}{normal SABR} (the
quotes' own model).
\textbf{5.} \omterm{def:rc:sabr-in-rates-and-the-volatility-cube:normal}{Normal SABR}, then the 3\% shift.
\textbf{6.} $0.177$, $0.280$ and $0.295$ basis points of annuity, times $2.0\times
500$ million: EUR~17\,717, 27\,999 and 29\,477.
\textbf{7.} EUR~10\,282, 37\% of the 3\% price.
\textbf{8.} The modelled rate cannot go below $-1\%$, and its volatility falls as the
rate nears the floor, so the probability of finishing below $-0.90\%$ is small.
\textbf{9.} The desk would rather sell at the highest defensible price, the client
buy at the lowest; the 1\% shift is the model most favourable to the client.
\textbf{10.} 0.90\% (exercise~7).
\textbf{11.} The difference between the two strikes' dynamics: the smile slope and
curvature below the quotes, i.e. the model's wing, plus the gap between the
strikes; the desk is exposed to the market re-marking the wing.
\textbf{12.} Rates near $-0.75\%$ would make the 1\% shift's floor almost binding and
its prices meaningless; the desk would need a larger shift, and marks would
jump when it changed.
\textbf{13.} Its choice moves the prices of unquoted strikes by tens of per cent
(here 37\%) while fitting quoted strikes equally well; that is model
uncertainty, and a reserve covers the spread of prices across admissible shifts.
\textbf{14.} Nothing: the rate is normal given the volatility, unbounded below,
which puts more weight on very negative rates than any shifted model.
\textbf{15.} Fit all candidate models; check densities over the book's strikes;
compare wing prices across models and with any traded points (dealer polls);
backtest the smile's behaviour when the forward moved.
\textbf{16.} To keep the floor well below any plausible rate, so that the wing of
the distribution, and the prices of low strikes, do not depend on where the
floor is; a small shift makes the smile hostage to it.
\textbf{17.} Reasonable within a currency if it is well below the \omterm{def:rc:sabr-in-rates-and-the-volatility-cube:elb}{effective lower
bound}; the choice is part of the market convention, which makes quotes
comparable.
\textbf{18.} That quoted strikes pin the price only near the forward, that three
standard models give EUR~17\,700 to 29\,500, and what the desk charges and why.
\textbf{19.} Named result: \emph{the minus-ninety strike} costs EUR~17\,717 with a 1\%
shift and EUR~27\,999 with 3\% (37\% apart), and 0.90\% is the smallest shift
that keeps the density positive from $-0.80\%$ up.
\textbf{20.} The model's assumption about how low rates can go and how volatile
they are near that bound.
\end{solution}

\begin{solution}{iq:rc:sabr-in-rates-and-the-volatility-cube:1}
By leaving lognormal (Black) quotes, which cannot represent negative forwards
or strikes: quoting normal (Bachelier) volatilities, or Black volatilities of a
shifted rate $F+\zeta$; models moved to \omterm{def:rc:sabr-in-rates-and-the-volatility-cube:normal}{normal SABR} or \omterm{def:rc:sabr-in-rates-and-the-volatility-cube:shifted}{shifted SABR}, later also
free-boundary SABR.
\iqlookfor{normal quotes and shifted models, and why Black fails.}
\end{solution}

\begin{solution}{iq:rc:sabr-in-rates-and-the-volatility-cube:2}
$\beta$ sets the backbone: how the at-the-money volatility moves with the
forward ($\beta=0$ flat in normal terms, $\beta=1$ flat in lognormal terms) and
part of the skew. Since $\rho$ also sets the skew, fitting both on one smile is
unstable; $\beta$ is fixed from the desk's view of the backbone or a historical
regression of at-the-money volatility on the forward.
\iqlookfor{backbone versus skew, and the identification problem.}
\end{solution}

\begin{solution}{iq:rc:sabr-in-rates-and-the-volatility-cube:3}
For each (expiry, tenor): forward and annuity from the multi-curve; clean the
quotes (normal vols or premiums, straddles and strangles); fix $\beta$ and the
shift; fit $\alpha,\rho,\nu$ with bounds; check errors and parameter smoothness
across the grid; interpolate parameters (variance in expiry); check density and
calendar arbitrage; publish with diagnostics.
\iqlookfor{a pipeline with checks, not only a fit.}
\end{solution}

\begin{solution}{iq:rc:sabr-in-rates-and-the-volatility-cube:4}
\omterm{def:rc:sabr-in-rates-and-the-volatility-cube:normal}{Normal SABR}: the smile shifts right by ten basis points, the at-the-money normal
volatility unchanged. Lognormal with $\beta=1$: the at-the-money Black volatility
is unchanged, so the normal volatility rises roughly in proportion to the
forward, and the smile moves in relative (log-moneyness) terms.
\iqlookfor{the two backbones.}
\end{solution}

\begin{solution}{iq:rc:sabr-in-rates-and-the-volatility-cube:5}
In the low-strike wing of long expiries (large $\nu^2T$), near the lower bound
of the rate, because the expansion is asymptotic and misrepresents the
absorbing boundary. Fixes: a larger shift, an exact or numerical SABR (PDE,
integration), free-boundary SABR, or wing extrapolation with a positive
density; always test the density over the book's strikes.
\iqlookfor{the cause and the fixes.}
\end{solution}

\begin{solution}{iq:rc:sabr-in-rates-and-the-volatility-cube:6}
Look at the inputs first: stale or crossed quotes, a forward from the wrong
curve, a strike grid shifted by the forward; then at the fit: starting point,
bounds hit ($\rho$ near $\pm1$, $\nu$ exploding), conversion failures of deep
quotes to the model's volatility; compare with yesterday's parameters and the
neighbouring sections; add a regularisation towards neighbours if the section
is thinly quoted.
\iqlookfor{data before optimiser, and a systematic check.}
\end{solution}
